% !TeX root = ../main.tex
\section{実験方法}
\label{sec:experiments}
\subsection{評価目的と共通設定}
提案モデルの構造選択と係数解釈を検証するため、Lorenz96で構造回復、線形VARで符号回復を評価する。系列長\(T\in\{250,500\}\)には同じ生成系列の先頭区間を用い、訓練区間で標準化する。評価する関係を変数間に限定し、自己辺は採点から除外する。

係数生成器は隠れ層1層のMLPとする。独立な予備3系列で幅\(\{8,32\}\)を比較し、各条件で選ばれた幅8を用いる。基準はLorenz96では平均F1、線形VARでは正・負のbalanced accuracy（BA）の平均である。構造回復では正則化10候補から平均F1が最大の候補を選び、同点なら平均選択辺数、尺度の小さい順に選ぶ。

正則化尺度\(s\)を\(\lambda_0=s/(Kp^2)\)へ換算する。係数生成器の学習率は\(0.003\)、ロジットの学習率は\(0.03\)、時間平滑化係数は\(0.01\)、重み減衰係数は\(0.001\)、初期選択確率は\(0.5\)とする。全標本をまとめて\(10{,}000\)エポック学習した最終モデルを評価する。Lorenz96では系列ごとに学習初期値を変え、線形VARでは固定する。

\subsection{Lorenz96における構造回復}
NGCで用いられるLorenz96系\citep{tank2022}
\begin{equation}
 \frac{dx_i(s)}{ds}
 =\left(x_{i+1}(s)-x_{i-2}(s)\right)x_{i-1}(s)-x_i(s)+F
 \label{eq:lorenz}
\end{equation}
から、周期境界条件の下で時系列を生成する。変数数\(p=20\)、外力\(F=40\)、観測間隔は約\(0.05\)とする。初期1000点を除き、標準偏差\(0.1\)の独立な正規観測雑音を加える。最大ラグ\(K=5\)とし、生成した5系列を評価する。データと学習初期値は手法間で対応させる。

参照構造は微分方程式の直接依存とし、自己辺を除くと\(i+1,i-1,i-2\)から\(i\)への60辺である。間接依存と観測雑音を含む離散系列から、この構造の回復を評価する。

GSE-VARの正則化尺度は\(\{0,0.3,1,2,3,4,5,7,10,30\}\)の10候補とする。予備評価で選んだ尺度は\(T=250,500\)の順に4、5である。NGCの採用値は0.5、0.4で、各系列長の10候補を独立に学習する。

\subsection{線形VAR過程における符号回復}
GVARの符号評価に用いられる4変量の生成法\citep{marcinkevics2021}に従い、\(\bm x_t=(x_t,w_t,y_t,z_t)^\top\)を
\begin{equation}
 \bm x_t=A\bm x_{t-1}+\bm\varepsilon_t,\quad
 A=\begin{pmatrix}
 a_1&0&0&0\\
 a_3&a_2&0&0\\
 0&a_5&a_4&0\\
 0&a_7&a_8&a_6
 \end{pmatrix},\quad
 \bm\varepsilon_t\sim N(0,0.4^2I)
 \label{eq:linear-generation}
\end{equation}
から生成する。各\(a_r\)は、等確率で選ぶ区間\([-0.8,-0.2]\)または\([0.2,0.8]\)の一様分布に従う。下三角行列の対角の絶対値が\(1\)未満なので過程は安定である。初期値\(0\)、最大ラグ\(K=1\)とし、係数を変えた10系列を評価する。非自己の真辺は4本であり、標準化後も正解の符号は変わらない。

正則化尺度は\(\{0,0.01,0.03,0.1,0.3,1,3,10,30\}\)の9候補とする。採用尺度は\(T=250,500\)の順に、GSE-VARが\(0.3\)、\(0.1\)である。

\subsection{比較手法と評価指標}
構造回復ではNGCのcMLP、GVAR、線形VARを比較対象とする。符号回復では、符号付き係数を持つGVARと線形VARを用いる。NGCは隠れ層1層・幅100、階層型グループ正則化を用い、5000エポック学習する。GVARは隠れ層2層・幅50、1000エポックとし、二値化には原系列と逆時間系列の推定構造の一致を用いる選択手順を適用する\citep{tank2022,marcinkevics2021}。線形VAR実験のGVARには、疎性係数\(0.2\)、平滑化係数\(0.5\)を用いる。Lorenz96のGVARは独立な予備系列で設定を選ぶ。

線形VARは切片付き最小二乗法で推定し、系列の全ラグを除く制約に対するF検定を行う。Benjamini--Hochberg法を水準\(0.05\)で適用し、二値グラフを得る\citep{benjamini1995}。

構造回復は、採用した設定で得られるグラフと、設定を変えた際の選択結果の両面から評価する。前者には適合率、再現率、F1と辺選択頻度を用い、予備探索で採用した正則化強度の結果を集計する。GSE-VARはゲートがいずれかのラグで非\(0\)となる辺、NGCは入力重みグループが非\(0\)となる辺を選択されたものとする。

後者にはPR曲線の面積（AUPRC）を用いる。GSE-VARとNGCでは正則化10候補を大から小へ変え、GVARでは\(S_{ij}=\max_k\median_t|c_{kij,t}|\)、線形VARでは負のF検定\(p\)値の閾値を変える。得られた適合率と再現率の隣接点間を台形則で積分するため、前2手法は正則化による構造変化、後2手法は単一モデルでの選択結果を評価する。選択辺がない場合の適合率は\(1\)と定義し、観測した再現率の範囲のみを積分する。

符号回復では、選択辺の係数の時点中央値の符号を用い、未選択辺は\(0\)とする。この符号から正辺・負辺をそれぞれ陽性とするBAを求め、その平均を主指標とする。BAは感度と特異度の平均であり、その他の変数対はそれぞれ陰性となる。片側のクラスがない系列は欠測とし、各指標を系列ごとの値、平均、標本標準偏差で示す。
